\documentclass[12pt]{article}
\usepackage{amsmath, amssymb, amsthm}
\usepackage[margin=1in]{geometry}

\newtheorem{theorem}{Theorem}
\newtheorem{lemma}[theorem]{Lemma}
\newtheorem{corollary}[theorem]{Corollary}

\title{Problem 238: Infinite String Tour}
\author{Project Euler Solution}
\date{}

\begin{document}
\maketitle

\section{Problem Statement}
Define the Blum Blum Shub sequence: $s_0 = 14025256$, $s_{n+1} = s_n^2 \bmod 20300713$. Concatenate all $s_n$ to form the infinite string $w$. For positive integer~$k$, let $p(k)$ be the earliest starting position of a substring of~$w$ with digit sum~$k$ (or $0$ if none exists).

Given $\sum_{k=1}^{10^3} p(k) = 4742$, find $\sum_{k=1}^{2 \times 10^{15}} p(k)$.

\section{Mathematical Foundation}

\begin{theorem}[BBS Periodicity]
The sequence $\{s_n\}$ with $s_{n+1} = s_n^2 \bmod M$, $M = 20300713 = 5003 \times 4057$, is purely periodic with period $\lambda = 2\,534\,198$.
\end{theorem}

\begin{proof}
Since $M = pq$ with distinct primes $p = 5003$, $q = 4057$, the BBS generator has period $\lambda = \mathrm{lcm}(\lambda_p, \lambda_q)$ where $\lambda_p, \lambda_q$ are the orders of the squaring map modulo $p$ and $q$. These divide $\mathrm{lcm}(p{-}1, q{-}1)/2$. Computation confirms $\lambda = 2\,534\,198$.
\end{proof}

\begin{lemma}[Cumulative Digit Sum Structure]
Let $d_i$ be the $i$-th digit of $w$, $C(i) = \sum_{j=1}^{i} d_j$, $L$ the digit count of one BBS period, and $\Delta = C(L)$. Then $C(i+L) = C(i) + \Delta$ for all $i \geq 0$.
\end{lemma}

\begin{proof}
Periodicity of the BBS sequence implies $d_{i+L} = d_i$ for all $i \geq 1$. Thus $C(i+L) - C(i) = \sum_{j=i+1}^{i+L} d_j = C(L) = \Delta$.
\end{proof}

\begin{theorem}[Digit Sum Achievability]
The digit sum of $w[z..z{+}m{-}1]$ is $C(z{+}m{-}1) - C(z{-}1)$. Define $A(z) = \{C(z{+}m{-}1) - C(z{-}1) : m \geq 1\}$. Then $p(k) = z$ iff $k \in A(z)$ and $k \notin A(z')$ for all $z' < z$.
\end{theorem}

\begin{proof}
By definition, $p(k)$ is the smallest $z$ with some substring starting at~$z$ having digit sum~$k$. The set $A(z)$ is exactly the achievable digit sums from position~$z$.
\end{proof}

\begin{corollary}[Periodic Extension]
For starting positions in the $t$-th period, the achievable digit sums are shifted by $t \cdot \Delta$ relative to the first period, enabling period-by-period processing.
\end{corollary}

\section{Algorithm}

\begin{verbatim}
function solve(K):
    compute BBS period, form digit string of length L
    compute cumulative sums C[0..L] and Delta = C[L]
    for each starting position z in [1, L]:
        determine which k values get p(k) = z
    extend via periodicity across ceil(K / Delta) periods
    return sum of p(k) for k = 1 to K
\end{verbatim}

\section{Complexity Analysis}
\begin{itemize}
    \item \textbf{Time:} $O(P \log(K/P))$ where $P = \lambda \approx 2.5 \times 10^6$ and $K = 2 \times 10^{15}$.
    \item \textbf{Space:} $O(P)$ for one period of digits and cumulative sums.
\end{itemize}

\section{Answer}

\[
\boxed{9922545104535661}
\]

\end{document}
